\section*{Abstract}
\thispagestyle{empty}
The Dynamic Request Acceptance and Unsplittable Scheduling Problem (DRAUSP) requires an online decision maker to accept or reject incoming requests and, if accepted, to schedule them as contiguous blocks across multiple capacity-limited resources, so as to maximize cumulated revenue over a finite horizon. \\
While exact stochastic dynamic programming solves DRAUSP optimally, it becomes computationally intractable as the number of resources and requests grows. This report formulates DRAUSP as a finite-horizon Markov Decision Process and learns an acceptance-and-scheduling policy via value-based deep reinforcement learning. A standard Deep Q-Network (DQN) with action masking and a reject-biased exploration scheme serves as the baseline. Since the problem exhibits domain structure that a generic network is not guaranteed to respect -- e.g., more remaining capacity or a higher-reward request should never yield a lower $Q$-value -- the goal of this work is to develop a Lattice-DQN that enforces this structure using a Deep Lattice Network architecture. We test the monotonicity of the learned $Q$-functions and compare the performance of both agents on a set of benchmark instances of \textcite{drausp_lion18} under a parameter-matched comparison. The Standard-DQN reaches $90$--$100\%$ of the optimal objective value on most instances and, without any structural constraints, already learns near-perfect monotonicity in most state components. The Lattice-DQN satisfies all monotonicity relations exactly by construction, but its more constrained function class leads to slower convergence and lower reward on harder, longer-horizon instances. These results suggest that structural monotonicity constraints guarantee reliable, interpretable $Q$-functions at the cost of reduced policy flexibility, motivating hybrid architectures as a direction for future work.

\clearpage
